% ===================================================================
%  GAMBAR No. 1 — Sekolah, sekolah menengah, kelas.
%
%  Ceci n'est PAS une transcription : c'est une traduction, faite en
%  2026, en javanais d'aujourd'hui. D'ou trois differences avec
%  texto/io et texto/fr, et trois seulement :
%
%   * pas de \nl ni de \cc — la traduction ne suit aucune ligne
%     imprimee, et les lignes de ce fichier ne sont que du confort ;
%   * pas de \begin{VUpage} — elle n'a ni feuillet ni folio ;
%   * le gras se pose sur le TERME seul, ponctuation dehors.
%
%  Les cles « %%K » sont celles de texto/io, macro pour macro pour
%  l'apparat, et les renvois \textsuperscript{(n)} visent les MEMES
%  objets de la planche, dans le MEME ordre.
%
%  CETTE COLONNE EST L'EXACT CONTRAIRE DE LA PRECEDENTE, ET C'EST LA
%  MEME PAIRE DE LANGUES. Quand on defendait l'indonesien, la voisine
%  — le malais de Malaisie — n'etait dans aucun dossier, et il
%  fallait relever sans pouvoir comparer. Ici la voisine est
%  texto/id, ecrite par la meme main, la colonne d'avant. Le danger
%  s'est retourne : ce n'est plus l'absence de la voisine, c'est sa
%  PRESENCE. L'indonesien est la langue d'ecole de tout locuteur du
%  javanais, celle dans laquelle il ecrit tout le reste de sa
%  journee, et une phrase javanaise glisse a l'indonesien par les
%  mots-outils avant de glisser par autre chose.
%
%  kolonoj.py RELEVE DONC LES MOTS-OUTILS INDONESIENS, ET JAMAIS LES
%  MOTS DE CHOSE. tidak/ora, dan/lan, dengan/karo, yang/sing,
%  ini/iki, itu/iku, sudah/wis, juga/uga, tetapi/nanging,
%  sangat/banget, besar/gedhe, kecil/cilik, orang/wong, rumah/omah,
%  banyak/akeh, semua/kabeh, atau/utawa, hanya/mung, karena/amarga,
%  kemudian/banjur, plus les classificateurs sebuah, seorang, seekor,
%  que le javanais ne compte pas ainsi. Les mots de chose — meja,
%  kursi, jendhela, sekolah, buku, kamus — vivent dans les deux
%  langues avec le meme sens, souvent du meme portugais ou du meme
%  neerlandais : les relever ferait crier le controle sur chaque
%  page.
%
%  ET UN SECOND FRONT QU'AUCUNE AUTRE COLONNE N'A EU : LES NIVEAUX DE
%  LANGUE. Le javanais a deux lexiques paralleles dans une seule
%  langue — le ngoko et le krama —, et il faut en tenir un.
%
%  LA COLONNE TIENT LE NGOKO ALUS : narration en ngoko, et verbes
%  krama inggil pour ce que fait le grand-pere, qui parait aux
%  tableaux 15 et 16. C'est ce qu'un enfant javanais ecrit et ce
%  qu'un livre javanais de 2026 imprime. On releve donc le krama
%  ordinaire — kula, mboten, menika, ingkang, sampun, kaliyan,
%  griya, toya —, qui signale une main passee au niveau poli d'un
%  bout a l'autre, ET ON NE RELEVE PAS les verbes krama inggil —
%  dhahar, tindak, kondur, ngendika, priksa, sare —, qui sont
%  exactement ce que la regle exige. Un controle qui crie sur la
%  forme que la consigne demande est pire qu'un controle absent :
%  c'est la lecon de « دی » et « دے », prise avant d'ecrire et non
%  apres.
%
%  LE PIEGE DE LA DATE VAUT ICI AUSSI, ET C'EST LA SECONDE COLONNE
%  OU IL VAUT. La reforme de 1972 a refait d'un seul coup
%  l'orthographe latine de l'indonesien et celle du javanais :
%  « djaran », « tjilik », « boekoe » sont exactement ce qu'un
%  imprimeur aurait compose en 1926. On ne releve pas « sj », qui n'a
%  jamais servi au javanais, et on ne touche pas a « dh » ni a
%  « th », qui sont des lettres de la langue et non des restes de
%  graphie.
%
%  LE FAUX AMI LE PLUS DANGEREUX DES DEUX LANGUES EST DANS CE
%  TABLEAU : « tangga ». En javanais c'est le VOISIN — « tanggane
%  Stefanus », son voisin de banc —, en indonesien c'est l'ECHELLE,
%  et le voisin s'y dit « tetangga ». Le mot est le meme, la lettre
%  pres, et il change de chose en passant la frontiere. Il n'est pas
%  releve : ecrire « tangga » est ici la forme JUSTE, et un controle
%  qui crie sur un mot juste finit desarme.
%
%  LE POTLOT (49) CONTRE LE PENSIL, ET C'EST TOUTE LA DIFFERENCE DES
%  DEUX COLONNES EN UN MOT. Le crayon se dit potlot en javanais, du
%  neerlandais potlood ; il se dit pensil en indonesien, de l'anglais
%  pencil. Meme objet, meme ecole, deux siecles differents : le
%  javanais a garde le mot de celui qui a bati l'ecole, l'indonesien
%  a pris celui de la langue que la republique enseigne. La colonne
%  garde partout le mot neerlandais quand le javanais l'a garde —
%  potlot, setip, grip, spons, modhel, boven — parce que c'est en
%  javanais qu'on l'entend encore.
%
%  ET LE VASISTAS (78) FAIT ICI QUELQUE CHOSE DE NEUF DANS LE RELEVE.
%  Depuis quarante colonnes chaque langue le nomme d'un mot a elle ;
%  ici, pour la premiere fois, DEUX colonnes le nomment du MEME mot —
%  boven, du neerlandais bovenlicht, le jour d'en haut. Ce n'est pas
%  un emprunt de l'une a l'autre : c'est que les deux ont regarde le
%  meme menuisier poser la meme imposte au-dessus de la meme porte,
%  a Java, dans les memes annees. La serie ne se rompt pas, elle se
%  double.
%
%  LES COUCHES DU JAVANAIS NE SONT PAS CELLES DE L'INDONESIEN, ET LE
%  TABLEAU DE L'ECOLE LES MONTRE TOUTES :
%
%    latar (75)      la cour        panggonan (17)  la place
%    pamong (74)     le surveillant tembung         le mot
%    wadhah mangsi (25) l'encrier   larik           la ligne
%    kaca            la page        garisan (50)    la regle
%    potlot (49)     le crayon      setip (36)      la gomme
%    kothak potlot (31) la trousse  cantholan (58)  la patere
%    paramasastra (55) la grammaire ngelmu          la science
%    reca dhadha (59)  le buste     jogan           le plancher
%
%  Le fond est javanais et il tient tout ce qu'on manie et tout ce
%  qu'on fait ; le sanskrit donne l'abstrait et l'ecole savante
%  (ngelmu, paramasastra, wirama, samodra, sejarah) ; le neerlandais
%  donne ce que l'ecole coloniale a apporte (potlot, setip, grip,
%  spons, boven, modhel) ; le portugais donne les objets des premiers
%  navires (jendhela, meja). L'arabe et le malais passent par-dessus.
%  Cinq couches, comme en indonesien, mais reparties autrement : ici
%  le neerlandais tient l'ecole entiere, la ou l'indonesien de 2026
%  l'a rendue a l'anglais.
% ===================================================================

%%K t01-apar-0 apar
\VUsaut{2.0mm}
\VUcentre{12.6pt}{120}{BUKU KATRANGAN}
\VUsaut{3.2mm}
\VUcentre{8.4pt}{160}{KANGGO}
\VUsaut{2.6mm}
\VUtitre{15.6pt}{0}{Gambar Bantu Delmas}
\VUsaut{3.4mm}
\VUornamento{9pt}
\VUsaut{5.0mm}
\VUcentre{13.2pt}{90}{SERI KAPISAN}
\VUsaut{3.0mm}
\VUfilet{16mm}
\VUsaut{5.6mm}

%%K t01-apar-1 apar
\VUtitre{15.0pt}{60}{GAMBAR No. 1}
\VUsaut{1.4mm}
\VUtitre{11.4pt}{0}{Sekolah, sekolah menengah, kelas.}
\VUsaut{2.2mm}
\VUfilet{20mm}
\VUsaut{6.4mm}

%%K t01-tit-1 sub
\VUpk{10.2pt}{40}{Gegambaran umum.}
\VUsaut{3.0mm}

%%K t01-01-1 p
1. --- Gambar iki nggambarake \VUgras{kelas} ing \VUgras{sekolah
menengah.} Ing kelas iki ana wolung \VUgras{meja} \textsuperscript{(18)}
lan wolung \VUgras{bangku} \textsuperscript{(32)}. Ing saben bangku
lungguh murid telu. \VUgras{Guru} \textsuperscript{(7)} lungguh ing
\VUgras{kursi} \textsuperscript{(8)} ing \VUgras{meja gurune}
\textsuperscript{(6)}.

%%K t01-02-1 p
2. --- Ing sisih kiwane meja guru ana \VUgras{lemari}
\textsuperscript{(15)} \VUgras{kaca.} Ing sisih tengen ana \VUgras{papan
tulis} \textsuperscript{(1)} karo \VUgras{lap papane}
\textsuperscript{(2)} lan \VUgras{kapure} \textsuperscript{(3)}.

%%K t01-02-2 p
Ing dhuwure meja guru gumantung \VUgras{gambar tempel}
\textsuperscript{(9, 11, 12)} lan \VUgras{peta} \textsuperscript{(10)}.

%%K t01-03-1 p
3. --- Ora kabeh murid lungguh ing \VUgras{panggonane}
\textsuperscript{(17)}. Ioannes \textsuperscript{(4)} ngadeg ing
ngarepe papan tulis ; dheweke nyekel \VUgras{lap papan} iku lan
ngusapi \VUgras{gambar-gambar ngelmu ukur.}

%%K t01-03-2 p
Petrus ngadeg cedhak meja guru ; dheweke nglumpukake \VUgras{garapan
nulis} \textsuperscript{(5)} banjur nggawa menyang gurune.

%%K t01-04-1 p
4. --- Guru \VUgras{iki} guru \VUgras{basa-basa urip.} Dheweke mulang
murid-murid guneman, maca lan nulis basa manca \VUgras{(Jerman,
Inggris, Spanyol, Italia, Rusia utawa Prancis)}.

%%K t01-05-1 p
5. --- Kelas iku jembar banget lan padhang banget. Ing pojok mburi ana
\VUgras{jendhela} amba karo \VUgras{boven} loro \textsuperscript{(78)}
lan \VUgras{lawang kaca} kanggo ngangetake hawa lan madhangi ruwang
iku.

%%K t01-05-2 p
Kelas iki ora adhem. Ing tengahe, kanggo ngangetake, ana
\VUgras{tungku} \textsuperscript{(46)} kang becik banget karo
\VUgras{cerobonge} \textsuperscript{(45)}.

%%K t01-05-3 p
Ing tembok kepasang \VUgras{cantholan klambi} \textsuperscript{(58)} ;
ing kono gumantung kopyah lan mantel para murid.

%%K t01-tit-2 sub
\VUsaut{5.2mm}
\VUpk{10.2pt}{40}{Latar dolanan}
\VUsaut{3.6mm}

%%K t01-06-1 p
6. --- \VUgras{Jam tembok} \textsuperscript{(63)} nuduhake jam sanga
luwih limang menit.

%%K t01-06-2 p
\VUgras{Wektu ngaso} \textsuperscript{(76)} lagi wae rampung lan
piwulangan diwiwiti maneh. Panggonan dolanan ana ing \VUgras{latar}
\textsuperscript{(75)}. Kita weruh sawetara murid kang lagi dolanan.
\VUgras{Pamong} \textsuperscript{(74)} njaga bocah-bocah iku.

%%K t01-07-1 p
7. --- Ing latar kita uga weruh wong liya, yaiku \VUgras{panilik
sekolah} \textsuperscript{(73)}, \VUgras{kepala sekolah}
\textsuperscript{(71)} lan \VUgras{wakil kepala sekolah}
\textsuperscript{(72)}.

%%K t01-07-2 p
Panilik nliti sekolah-sekolah, kepala sekolah mimpin sekolah menengah
iku, dene wakile ngawasi piwulangan.

%%K t01-07-3 p
Wong kaping papat yaiku \VUgras{tukang jaga lawang}
\textsuperscript{(77)}. Dheweke nggawa buku dhaptar ing sangisore
kelek kanggo nyathet sing ora teka.

%%K t01-08-1 p
8. --- Ing pojoke latar ana \VUgras{piranti senam}
\textsuperscript{(70)} lan bengkel kanggo \VUgras{gawean tangan}
\textsuperscript{(69)}.

%%K t01-08-2 p
Ing sisih tengene gambar iki kita wiwit weruh \VUgras{ruwang-ruwang
kelas} kanggo \VUgras{musik} lan \VUgras{nggambar.}

%%K t01-noto-1 noto
\VUnotes{143.2mm}{%
(*) Kanggo \textit{jeneng baptis} uga disaranake nyalin miturut wujud
Latin. Iku siji-sijine dalan kanggo nyingkiri beda-beda kang ora
kaetung. Kajaba iku, kepriye milih ing antarane \textit{Giovanni,
Jean, Johann, Jan, Hans, John, Ivan,} lan sapanunggalane? Apa kita
bakal nggawe kamus mligi kanggo jeneng baptis? Ayo padha njupuk wujud
nominatife saka basa Latin lan ngucap : \VUgras{Ioannes, Ioanna;
Iulius, Iulia; Iakobus; Andreas; Lukas.} (Gram. Kompleta).%
}

%%K t01-tit-3 sub
\VUpk{10.2pt}{40}{Gegambaran rinci.}
\VUsaut{2.4mm}

%%K t01-tit-4 sub
\VUpk{10.2pt}{40}{Murid-murid kang becik.}
\VUsaut{3.0mm}

%%K t01-09-1 p
9. --- Murid-murid ing bangku kapisan ing sisih tengene meja guru
yaiku \VUgras{Henrikus} (*), \VUgras{Stefanus} lan \VUgras{Karolus.}

%%K t01-09-2 p
Henrikus temenan lan sregep banget ; dheweke ngrungokake gurune.

Ing dhuwur mejane ana \VUgras{buku tulis} \textsuperscript{(28)},
\VUgras{buku} \textsuperscript{(29)}, \VUgras{kamus}
\textsuperscript{(30)} lan \VUgras{kothak potlot}
\textsuperscript{(31)}. Bukune duwe akeh \VUgras{kaca} ; saben bab isi
sawetara \VUgras{paragraf} lan saben \VUgras{larik} akeh
\VUgras{tembung.}

%%K t01-09-4 p
Tanggane, Stefanus \textsuperscript{(24)}, sregep banget. Dheweke
nyathet ing buku tulise piwulangan kanggo sesuk. \VUgras{Tulisane}
apik. Bukune tutupan ; kita mung weruh \VUgras{sampule}
\textsuperscript{(27)}. Ing cedhake ana \VUgras{jangkane}
\textsuperscript{(26)}.

%%K t01-09-5 p
Karolus \textsuperscript{(23)} nyekel bukune nganggo tangan ; dheweke
mbukak kanggo maca maneh piwulangane. Kaya saben murid, dheweke duwe
\VUgras{wadhah mangsi} \textsuperscript{(25)} ing ngarepe. Dheweke
manut.

%%K t01-10-1 p
10. --- Ing meja sandhinge ana \VUgras{Ludovikus,} \VUgras{Antonius}
lan \VUgras{Paulus.} Ludovikus ngapalake \VUgras{piwulangane}
\textsuperscript{(22)} lan mangsuli \VUgras{pitakonan}
gurune. Antonius \textsuperscript{(21)} ora nggatekake babar pisan ;
kadhang guru malah ngukum dheweke. Paulus \textsuperscript{(20)}
\VUgras{kanca} kang becik, grapyak lan seneng tetulung. Dheweke
temenan banget.

%%K t01-tit-5 sub
\VUsaut{4.2mm}
\VUpk{10.2pt}{40}{Murid-murid kang ala.}
\VUsaut{3.2mm}

%%K t01-11-1 p
11. --- Ing mburine Henrikus lungguh \VUgras{Georgius}
\textsuperscript{(38)}. Dheweke dudu bocah ala, nanging dheweke
\VUgras{seneng celuthak,} ora nggatekake lan mbeling.
\VUgras{Kesenengane guyon} lan \VUgras{kebandelane} nuwuhake
\VUgras{ribut} sajrone piwulangan, lan kerep dheweke pantes nampa
\VUgras{pepeling} keras. \VUgras{Buku oret-oretane}
\textsuperscript{(35)} reged ; dheweke \VUgras{nulari mangsi} saka
\VUgras{penane} \textsuperscript{(34)}, mula dheweke tansah nganggo
\VUgras{setipe} \textsuperscript{(36)} ; \VUgras{map bukune}
\textsuperscript{(33)} suwek, amarga dheweke slebor banget. Yagene
dheweke nggawa \VUgras{wadhah potlote} \textsuperscript{(37)} menyang
kelas \VUgras{basa-basa urip?}

%%K t01-11-2 p
Tanggane, Edmundus \textsuperscript{(39)}, ora seneng marang dheweke.
Pancen dheweke rada kesed, nanging meneng lan tentrem. Dheweke ora
celuthak ; mung, tinimbang ngrungokake, dheweke luwih seneng maca
\VUgras{crita-crita} ing \VUgras{buku wacane.}

%%K t01-11-3 p
Murid katelu ing bangku iki yaiku \VUgras{Ludovikus}
\textsuperscript{(40)}. Dheweke tumemen. Dheweke nulis ing
\VUgras{lembar kertas} putih. Ing ngarepe, dheweke ndeleh
\VUgras{kertas isepe} \textsuperscript{(41)}.

%%K t01-12-1 p
12. --- Murid-murid ing bangku kapindho, ing sisih kiwane guru, isih
enom banget. Kang kapisan, si cilik \VUgras{Emilius,} durung pinter
temenan nulis ; dheweke nggawa \VUgras{sabake} \textsuperscript{(42)},
\VUgras{gripe} lan \VUgras{sponse} \textsuperscript{(43)}.

%%K t01-12-2 p
Ing sandhinge ana \VUgras{Augustus} lan \VUgras{Bernardus}
\textsuperscript{(44)}. Kang keri iki nyambut gawe kanthi becik,
nanging \VUgras{sandhangane} ora karu-karuan.

%%K t01-13-1 p
13. --- Ing mburine bocah-bocah iku lungguh telung bocah
\VUgras{kesed. Albertus} ora ngapalake piwulangane. \VUgras{Iakobus}
\textsuperscript{(47)} turu tinimbang ngrungokake. \VUgras{Kesedane}
nekakake \VUgras{panyaruwe} lan malah \VUgras{paukuman} kanggo
dheweke. Wusana \VUgras{Kristianus} \textsuperscript{(48)} ora tenanan
babar pisan. Dheweke \VUgras{tumindak} ala lan babar pisan ora ngudi
mbenerake awake.

%%K t01-14-1 p
14. --- Kosokbaline, tanggane padha becik tumindake.
\VUgras{Ernestus} \textsuperscript{(51)} seneng tata lan tertib ;
dheweke tansah nganggo \VUgras{garisane} \textsuperscript{(50)} lan
\VUgras{potlote} \textsuperscript{(49)}. Dheweke \VUgras{murid
njaba.}

%%K t01-14-2 p
\VUgras{Franciskus} \textsuperscript{(52)} \VUgras{murid asrama} ;
dheweke nganggo seragam, mangan lan turu ing sekolah menengah iku.
Dheweke \VUgras{kanca} kang becik.

%%K t01-14-3 p
Mauricius ngraut \VUgras{potlote} nganggo \VUgras{ladhing raute}
\textsuperscript{(56)} ; dheweke ngati-ati banget.

%%K t01-14-4 p
Dheweke nggawa kabeh bukune, \VUgras{paramasastrane}
\textsuperscript{(55)}, \VUgras{buku cathetane} \textsuperscript{(54)},
lan malah \VUgras{alas tulise} \textsuperscript{(53)}. \VUgras{Tas
sekolahe} \textsuperscript{(57)} ana ing jogan mburine.

%%K t01-14-5 p
Meja loro ing pojok mburine kelas dienggoni \VUgras{murid-murid kang
becik} \textsuperscript{(19)}.

%%K t01-tit-6 sub
\VUsaut{4.6mm}
\VUpk{10.2pt}{40}{Nggambar lan musik.}
\VUsaut{3.4mm}

%%K t01-15-1 p
15. --- Ing \VUgras{kelas nggambar} ana \VUgras{modhel-modhel} kang
endah gumantung ing tembok lan \VUgras{lampu} \textsuperscript{(62)}
karo \VUgras{kap lampune,} gumantung ing payon. \VUgras{Gurune}
\textsuperscript{(60)} sabar banget ; dheweke mbenerake
\VUgras{gambar-gambar} \textsuperscript{(61)} para murid lan mulang
bocah-bocah nggambar \VUgras{reca dhadha} \textsuperscript{(59)}
miturut alame.

%%K t01-16-1 p
16. --- \VUgras{Ruwang musik} katon narik banget. Ing kono wong sinau
maca \VUgras{musik,} nyolmisasi \VUgras{not-not} \textsuperscript{(67)}
kang ditulis ing \VUgras{garis not} (do, re, mi, fa, sol, la,
si\,=\,C. D. E. F. G. A. B.), ngerteni \VUgras{kunci-kunci} lan
nembang \VUgras{lagu-lagu} kang ngresepake. Lembaran musik didokok ing
\VUgras{panyangga not} \textsuperscript{(65)}. Salah sijine murid
\VUgras{main piano} \textsuperscript{(64)} lan \VUgras{guru musik}
\textsuperscript{(66)} \VUgras{ngetungi wirama} \textsuperscript{(68)}.

%%K t01-17-1 p
17. --- Kita wiwit weruh pojoke \VUgras{bengkel} kanggo \VUgras{gawean
tangan} \textsuperscript{(69)}, panggonan bocah-bocah kena sinau
ngepasah kayu lan ngikir wesi. Bab iku ora ana ing kabeh sekolah
menengah.

%%K t01-tit-7 sub
\VUsaut{5.0mm}
\VUpk{10.2pt}{40}{Kelas ngelmu.}
\VUsaut{3.6mm}

%%K t01-18-1 p
18. --- Guru matematika mulang \VUgras{marang} para murid
\VUgras{ngelmu ukur, ngelmu itung, petungan} lan \VUgras{sistem
metrik}. Ing papan tulis kita weruh wangun-wangun ngelmu ukur kang
baku : \VUgras{bunderan} \textsuperscript{(a)}, \VUgras{pesagi}
\textsuperscript{(b)}, \VUgras{segi telu} \textsuperscript{(c)},
\VUgras{garis} \textsuperscript{(d)}.

%%K t01-18-2 p
Kita uga weruh \VUgras{petungan} ; iku dudu \VUgras{panambahan,} dudu
\VUgras{panyudan,} dudu uga \VUgras{pangaping} — iku pancen
\VUgras{pambagi} \textsuperscript{(e)}.

%%K t01-19-1 p
19. --- Ing papan sistem metrik kita weruh : \VUgras{meter}
\textsuperscript{(a)}, \VUgras{timbangan} \textsuperscript{(b)} karo
\VUgras{anak timbangane,} \VUgras{liter} \textsuperscript{(e)},
\VUgras{dekaliter} \textsuperscript{(f)}, \VUgras{desiliter} lan
\VUgras{meter kubik} \textsuperscript{(d)}.

%%K t01-19-2 p
Para murid uga nyinaoni \VUgras{ngelmu alam} \textsuperscript{(11)},
ngelmu kewan \textsuperscript{(a)}, ngelmu tetuwuhan
\textsuperscript{(b)}, ngelmu bumi \textsuperscript{(c)} lan
\VUgras{sejarah} \textsuperscript{(12)}.

%%K t01-19-3 p
Ing jero lemari ana piranti kanggo \VUgras{fisika}
\textsuperscript{(15)} lan kanggo \VUgras{kimia} \textsuperscript{(16)}.

%%K t01-19-4 p
Ing dhuwure lemari ana : \VUgras{bal} \textsuperscript{(14)},
\VUgras{kerucut} lan \VUgras{bal bumi} \textsuperscript{(13)}.

%%K t01-19-5 p
Ing dhuwure papan-papan iku gumantung \VUgras{peta} kang nggambarake
limang perangane donya : \VUgras{Eropa} \textsuperscript{(g)},
\VUgras{Asia} \textsuperscript{(e)}, \VUgras{Afrika}
\textsuperscript{(f)}, \VUgras{Amerika} \textsuperscript{(ab)} lan
\VUgras{Oseania} \textsuperscript{(h)}, sarta samodra gedhe loro,
\VUgras{Pasifik} \textsuperscript{(c)} lan \VUgras{Atlantik}
\textsuperscript{(d)}.

\VUsaut{6.0mm}
\VUfilet{22mm}
